\renewcommand{\arraystretch}{1.15}

\section{Experimental Setup Details} \label{app:setup}

\paragraph{Datasets and splits.}
We evaluate on four domains presented sequentially: 
GSM8K \citep{cobbe_training_2021} (\texttt{train} split), 
MMLU \citep{hendrycks_measuring_2021} (\texttt{test} split), 
ARC-Challenge \citep{clark_think_2018} (\texttt{test} split), 
and TruthfulQA \citep{lin_truthfulqa_2021} (\texttt{validation} split, MC1 variant).
We sample 500 questions per domain for a total of 2{,}000 questions per run.
The default domain order is GSM8K $\to$ MMLU $\to$ ARC $\to$ TruthfulQA.

\paragraph{Models.}
We evaluate four instruction-tuned small language models spanning 2--8B parameters:
Llama~3.2-3B-Instruct and Llama~3.1-8B-Instruct \citep{grattafiori_llama_2024},
Gemma~2-2B-IT \citep{team_gemma_2024},
and Phi~3.5-Mini-Instruct (3.8B) \citep{abdin_phi3_2024}.

\paragraph{Trainable parameters.}
Table~\ref{tab:lora_params} reports the exact LoRA parameter counts per model.
The variation reflects architectural differences (Phi uses a fused \texttt{qkv\_proj} 
module rather than separate \texttt{q\_proj}/\texttt{v\_proj}).

\begin{table}[htbp]
\centering
\small
\resizebox{\columnwidth}{!}{%
\begin{tabular}{lcccc}
\toprule
\textbf{Model} & \textbf{Total params} & \textbf{LoRA layers} & \textbf{LoRA params} & \textbf{\% of total} \\
\midrule
Llama 3.2-3B  & 3.21B & 4  & 327{,}680 & 0.010\% \\
Llama 3.1-8B  & 8.03B & 8  & 1{,}048{,}576 & 0.013\% \\
Gemma 2-2B    & 2.61B & 8  & 491{,}520 & 0.019\% \\
Phi 3.5-Mini  & 3.82B & 8  & 786{,}432 & 0.021\% \\
\bottomrule
\end{tabular}}%
\caption{LoRA trainable parameters per model (rank $r{=}8$).
Llama~3.2-3B: last 4 layers (24--27 of 28);
Llama~3.1-8B: last 8 layers;
Gemma and Phi: last 8 layers
(Gemma: 18--25 of 26; Phi: 24--31 of 32).}
\label{tab:lora_params}
\end{table}

\paragraph{Compute and software.}
All experiments were conducted on a shared workstation equipped with NVIDIA A100 (80\,GB) and RTX A6000 (48\,GB) GPUs.
Each 2{,}000-question SECL run takes approximately 2--4 GPU-hours on a single GPU, depending on model size.
We use PyTorch~2.9 \citep{paszke_pytorch_2019}, HuggingFace Transformers~4.57 \citep{wolf_transformers_2020}, and the PEFT library (v0.17) with CUDA~12.8.

\section{Metric Definitions} \label{app:metrics}

\emph{Expected Calibration Error} (ECE; \citealt{naeini_obtaining_2015, guo_calibration_2017}; lower is better). Predictions are divided into $M{=}10$ equal-width bins $B_1, \dots, B_M$ by confidence value:
\begin{equation}
    \text{ECE} = \sum_{m=1}^{M} \frac{|B_m|}{N} 
    \,\big|\text{acc}(B_m) - \text{conf}(B_m)\big|,
\end{equation}
where $N$ is the total number of predictions, $\text{acc}(B_m)$ is the fraction of correct predictions in bin $B_m$, and $\text{conf}(B_m)$ is the mean predicted confidence. A perfectly calibrated model achieves $\text{ECE} = 0$.

\emph{Brier score} \cite{brier_verification_1950} (lower is better). The mean squared error between predicted confidence and binary correctness:
\begin{equation}
    \text{BS} = \frac{1}{N} \sum_{i=1}^{N} 
    (c_i - y_i)^2,
\end{equation}
where $c_i \in [0,1]$ is the predicted confidence and $y_i \in \{0,1\}$ indicates correctness. The Brier score decomposes into calibration and refinement components.

\emph{AUROC} (area under the ROC curve; higher is better). The area under the receiver operating characteristic curve, treating predicted confidence as a score for binary classification of correctness. AUROC measures discrimination, whether higher confidence corresponds to higher correctness, independent of calibration.

\emph{Adaptive ECE} (AdaECE). Same as ECE but with equal-mass bins (lower is better); we report it to verify robustness to binning choice.

\emph{Accuracy}. Fraction of correct answers; we report it to verify that calibration updates do not degrade task performance.

\section{Hyperparameter Settings} \label{app:hyperparams}

\begin{table}[H]
\centering
\small
\resizebox{\columnwidth}{!}{%
\begin{tabular}{llc}
\toprule
\textbf{Component} & \textbf{Parameter} & \textbf{Value} \\
\midrule
\multirow{6}{*}{LoRA}
  & Rank $r$ & 8 \\
  & Scaling factor $\alpha$ & 16 \\
  & Target layers (Llama 3.2-3B) & Last 4 \\
  & Target layers (Llama 3.1-8B, Gemma, Phi) & Last 8 \\
  & Target modules (Llama, Gemma) & \texttt{q\_proj}, \texttt{v\_proj} \\
  & Target modules (Phi) & \texttt{qkv\_proj} \\
\midrule
\multirow{3}{*}{Optimization}
  & Optimizer & AdamW \\
  & Learning rate & $5 \times 10^{-5}$ \\
  & Epochs per question & 3 \\
\midrule
\multirow{2}{*}{Directional loss}
  & Step size $\alpha_{\text{step}}$ & 0.5 \\
  & Clip bound $\delta$ & 0.15 \\
\midrule
\multirow{1}{*}{Bin-gate}
  & Threshold & 1 bin \\
\midrule
\multirow{4}{*}{Page-Hinkley}
  & Tolerance $\epsilon$ & 0.05 \\
  & Detection threshold $\lambda$ & 3.0 \\
  & EMA smoothing $\alpha_{\text{ema}}$ & 0.05 \\
  & Warmup period & 30 questions \\
\midrule
\multirow{1}{*}{Burst}
  & Burst size $B$ & 50 \\
\midrule
\multirow{4}{*}{Normalization $\tau$}
  & Llama 3.2-3B & \textbf{0.7} \\
  & Gemma 2-2B & \textbf{1.5} \\
  & Phi 3.5-Mini (3.8B) & \textbf{1.5} \\
  & Llama 3.1-8B & \textbf{3.0} \\
\midrule
\multirow{2}{*}{Other}
  & Weight accumulation & On \\
  & Distractors $k$ & 4 \\
\bottomrule
\end{tabular}}%
\caption{Full hyperparameter settings used across all experiments.
Model-specific values are noted where applicable.}
\label{tab:hyperparams}
\end{table}

\section{Main Results} \label{app:main_results}

Table~\ref{tab:main_results_full} reports overall metrics for all methods and all four models.
We report both standard ECE (10 equal-width bins) and Adaptive ECE (AdaECE; 10 equal-mass bins; see Section~\ref{app:metrics}) to verify that improvements are robust to the choice of binning strategy.
Table~\ref{tab:main_results} in the main paper summarizes the same four models with ECE/Brier/AUROC and cost; this table adds AdaECE and supervised baselines.

\begin{table*}[htbp]
\centering
\footnotesize
\begin{tabular}{llccccc}
\toprule
\textbf{Model} & \textbf{Method}
  & \textbf{ECE}$\downarrow$ & \textbf{AdaECE}$\downarrow$ & \textbf{Brier}$\downarrow$
  & \textbf{AUROC}$\uparrow$ & \textbf{Acc} \\
\midrule
\multirow{6}{*}{Llama 3.2-3B}
  & Verbalized     & .170 & .167 & .292 & .510 & .576 \\
  & Self-Consistency$^\dagger$  & .093 & .096 & \textbf{.211} & \textbf{.728} & .624 \\
  & ~~+ Temp Scaling$^*$ & .047 & .075 & .250 & .504 & .576 \\
  & P(True) Norm        & .065 & .067 & .223 & .694 & .576 \\
  & ~~+ Temp Scaling$^*$ & \textbf{.029} & \textbf{.030} & .218 & .692 & .576 \\
  & SECL (Ours)         & \underline{.050} & \underline{.060} & .241 & .587 & .577 \\
\midrule
\multirow{5}{*}{Gemma 2-2B}
  & Verbalized     & .256 & .252 & .314 & .558 & .516 \\
  & ~~+ Temp Scaling$^*$ & .047 & .037 & .249 & .550 & .516 \\
  & P(True) Norm        & .141 & .142 & .259 & \textbf{.650} & .516 \\
  & ~~+ Temp Scaling$^*$ & \textbf{.037} & \textbf{.033} & \textbf{.233} & .647 & .516 \\
  & SECL (Ours)         & \underline{.056} & \underline{.060} & .254 & .548 & .515 \\
\midrule
\multirow{5}{*}{Phi 3.5-Mini (3.8B)}
  & Verbalized     & .251 & .240 & .275 & .600 & .667 \\
  & ~~+ Temp Scaling$^*$ & \textbf{.047} & \textbf{.049} & .215 & .583 & .667 \\
  & P(True) Norm        & .154 & .151 & .227 & \textbf{.675} & .667 \\
  & ~~+ Temp Scaling$^*$ & .059 & .060 & \textbf{.205} & .664 & .667 \\
  & SECL (Ours)         & \underline{.110} & \underline{.119} & .251 & .521 & .665 \\
\midrule
\multirow{5}{*}{Llama 3.1-8B}
  & Verbalized     & .225 & .225 & .258 & .684 & .644 \\
  & ~~+ Temp Scaling$^*$ & \textbf{.080} & .089 & .213 & .681 & .644 \\
  & P(True) Norm        & .120 & .117 & .211 & \textbf{.718} & .646 \\
  & ~~+ Temp Scaling$^*$ & .108 & .103 & \textbf{.208} & .716 & .646 \\
  & SECL (Ours)         & \underline{.083} & \textbf{.080} & .222 & .643 & .646 \\
\bottomrule
\end{tabular}
\vspace{0.3em}
\par\noindent{\footnotesize $^*$Requires ground-truth labels; fitted via 5-fold cross-validation. See Section~\ref{app:posthoc}.}
\par\noindent{\footnotesize $^\dagger$Self-Consistency requires $N{=}10$ sampling
passes per question. We report it for Llama~3.2-3B only to
contextualize SECL's single-pass efficiency; it is omitted for
other models as it measures generation consistency rather than
calibration directly.}
\caption{Overall results across all models and methods (2{,}000 questions each).
\textbf{Bold}: best overall per model for ECE, AdaECE, Brier, and AUROC.
\underline{Underline}: best among label-free methods.}
\label{tab:main_results_full}
\end{table*}

\subsection{Llama Multi-Seed Robustness} \label{app:llama_multiseed}

To quantify seed sensitivity for the main SECL setting, we report Llama~3.2-3B results for seeds 42, 43, and 44 on the same 2{,}000-question sequential stream for both the verbalized baseline and SECL.

\begin{table}[htbp]
\centering
\small
\resizebox{\columnwidth}{!}{%
\begin{tabular}{lcccc}
\toprule
\textbf{Metric} & \textbf{Seed 42} & \textbf{Seed 43} & \textbf{Seed 44} & \textbf{Mean$\pm$Std} \\
\midrule
Verbalized ECE & .1695 & .1757 & .1757 & .1736$\pm$.0029 \\
SECL ECE & .0496 & .0502 & .0392 & \textbf{.0464$\pm$.0051} \\
\midrule
Verbalized AUROC & .5095 & .5119 & .5119 & .5111$\pm$.0011 \\
SECL AUROC & .5870 & .6005 & .5920 & \textbf{.5932$\pm$.0056} \\
\bottomrule
\end{tabular}
}
\caption{Llama~3.2-3B multi-seed robustness (seeds 42, 43, 44). SECL remains strongly better calibrated than the verbalized baseline across seeds and also improves AUROC on this model.}
\label{tab:llama_multiseed}
\end{table}

\section{Post-Hoc Calibration Baselines} \label{app:posthoc}

We compare SECL against standard post-hoc calibration methods:
\textbf{temperature scaling} \citep{guo_calibration_2017} (single parameter $T$
fitted to minimize NLL) and \textbf{Platt scaling} \citep{platt_probabilistic_1999}
(logistic regression $\sigma(a \cdot c + b)$).
Both are applied to the soft verbalized confidence and fitted via 5-fold 
cross-validation (each fold uses 1{,}600 labeled examples for fitting).

\paragraph{Key difference from SECL.}
Temperature and Platt scaling are \emph{supervised} post-hoc methods:
they require ground-truth correctness labels to fit their parameters.
SECL is \emph{unsupervised} at test time; it uses only the model's own 
P(True) signal, with no access to ground truth.
Post-hoc methods are also monotone transforms that 
\emph{cannot} improve discrimination (AUROC);
SECL can improve AUROC by adapting the model's representations.

\begin{table*}[htbp]
\centering
\footnotesize
\begin{tabular}{llccccc}
\toprule
\textbf{Model} & \textbf{Method}
  & \textbf{ECE}$\downarrow$ & \textbf{AdaECE}$\downarrow$
  & \textbf{Brier}$\downarrow$ & \textbf{AUROC}$\uparrow$ 
  & \textbf{Conf.\ range} \\
\midrule
\multirow{4}{*}{Llama 3.2-3B}
  & Verbalized + Temp ($T{=}17.6$)  & .047 & .075 & .250 & .504 & [.44,\,.56] \\
  & Verbalized + Platt              & .021 & .057 & .244 & .481 & [.56,\,.61] \\
  & P(True) Norm + Temp ($T{=}1.8$)  & .029 & .030 & .218 & .692 & [.14,\,.86] \\
  & SECL (no labels)        & .050 & .060 & .241 & .587 & [.05,\,.85] \\
\midrule
\multirow{4}{*}{Gemma 2-2B}
  & Verbalized + Temp ($T{=}11.2$)  & .047 & .037 & .249 & .550 & [.41,\,.59] \\
  & Verbalized + Platt              & .035 & .050 & .250 & .550 & [.42,\,.55] \\
  & P(True) Norm + Temp ($T{=}3.0$)  & .037 & .033 & .233 & .647 & [.26,\,.74] \\
  & SECL (no labels)        & .056 & .060 & .254 & .548 & [.05,\,.95] \\
\midrule
\multirow{4}{*}{Phi 3.5-Mini (3.8B)}
  & Verbalized + Temp ($T{=}3.4$)   & .047 & .049 & .215 & .583 & [.29,\,.71] \\
  & Verbalized + Platt              & .052 & .050 & .216 & .585 & [.16,\,.69] \\
  & P(True) Norm + Temp ($T{=}2.4$)  & .059 & .060 & .205 & .664 & [.21,\,.79] \\
  & SECL (no labels)        & .110 & .119 & .251 & .521 & [.05,\,.95] \\
\midrule
\multirow{4}{*}{Llama 3.1-8B}
  & Verbalized + Temp ($T{=}2.8$)   & .080 & .089 & .213 & .681 & [.35,\,.74] \\
  & Verbalized + Platt              & .063 & .052 & .206 & .677 & [.02,\,.75] \\
  & P(True) Norm + Temp ($T{=}0.7$)  & .108 & .103 & .208 & .716 & [.02,\,.99] \\
  & SECL (no labels)        & .083 & .080 & .222 & .643 & [.25,\,.95] \\
\bottomrule
\end{tabular}
\caption{Post-hoc calibration baselines (5-fold CV). 
Temperature scaling compresses the confidence range substantially,
achieving low ECE by predicting near-constant values for models with
weak verbalized discrimination (fitted $T$ shown).}
\label{tab:posthoc}
\end{table*}

\paragraph{Discussion.}
For Llama, temperature scaling on the soft verbalized confidence requires 
$T{=}17.6$, compressing all predictions to [0.44,\,0.56], essentially 
predicting the base rate.
Platt scaling is more extreme: its output range [0.56,\,0.61] is near-constant,
yielding ECE$=$0.021 but \emph{degrading} AUROC to 0.481 (below chance).
These methods achieve low ECE by destroying discriminative information.
By contrast, SECL achieves ECE$=$0.050 while \emph{improving} AUROC
from 0.510 to 0.587 (+0.077), i.e., genuine calibration rather
than confidence compression.

For Phi, where the verbalized confidence carries more signal 
(baseline AUROC$=$0.600), temperature scaling uses a more moderate 
$T{=}3.4$ and achieves good calibration (ECE$=$0.047).
The strongest supervised baseline across all models is P(True) Norm~+~Temp,
which combines the discriminative P(True) signal with post-hoc recalibration.

\paragraph{Complementarity: SECL + Temperature Scaling.}
SECL and post-hoc calibration are complementary: applying temperature 
scaling to SECL's output further reduces ECE, since SECL improves the 
underlying confidence signal while temperature scaling recalibrates it.

\begin{table}[htbp]
\centering
\small
\resizebox{\columnwidth}{!}{%
\begin{tabular}{llcccc}
\toprule
\textbf{Model} & \textbf{Method}
  & \textbf{ECE}$\downarrow$ & \textbf{AdaECE}$\downarrow$
  & \textbf{Brier}$\downarrow$ & \textbf{AUROC}$\uparrow$ \\
\midrule
\multirow{3}{*}{Llama 3.2-3B}
  & P(True) Norm + Temp$^*$  & .029 & .030 & .218 & .692 \\
  & SECL            & .050 & .060 & .241 & .587 \\
  & SECL + Temp$^*$ & .049 & .053 & .241 & .584 \\
\midrule
\multirow{3}{*}{Gemma 2-2B}
  & P(True) Norm + Temp$^*$  & .037 & .033 & .233 & .647 \\
  & SECL            & .056 & .060 & .254 & .548 \\
  & SECL + Temp$^*$ & .011 & .037 & .249 & .542 \\
\midrule
\multirow{3}{*}{Phi 3.5-Mini (3.8B)}
  & P(True) Norm + Temp$^*$  & .059 & .060 & .205 & .664 \\
  & SECL            & .110 & .119 & .251 & .521 \\
  & SECL + Temp$^*$ & .097 & .082 & .232 & .516 \\
\midrule
\multirow{3}{*}{Llama 3.1-8B}
  & P(True) Norm + Temp$^*$  & .108 & .103 & .208 & .716 \\
  & SECL            & .083 & .080 & .222 & .643 \\
  & SECL + Temp$^*$ & .062 & .045 & .213 & .638 \\
\bottomrule
\end{tabular}}%
\vspace{0.3em}
\newline
{\small $^*$Requires ground-truth labels for temperature fitting.}
\caption{SECL combined with temperature scaling (5-fold CV).
SECL+Temp achieves the best or near-best ECE overall while preserving
SECL's discrimination improvements.}
\label{tab:secl_plus_temp}
\end{table}

\paragraph{Calibration--discrimination trade-off.}
SECL's primary objective is calibration (ECE), not discrimination (AUROC).
Table~\ref{tab:main_results_full} reveals a model-dependent trade-off:
for Llama, SECL \emph{improves} AUROC from .510 to .587 (+.077)
while reducing ECE from .170 to .050,
i.e., the adapted confidence carries better signal.
For Gemma, AUROC is roughly preserved (.558 $\to$ .548, $-$.010)
despite a large ECE reduction (.256 $\to$ .056).
For Phi, the trade-off is more pronounced:
AUROC drops from .600 to .521 ($-$.079) as ECE improves from .251 to .110.
Brier score, which jointly captures calibration and
discrimination, improves for all four models
(Llama~3B: .292 $\to$ .241; Llama~8B: .258 $\to$ .222; Gemma: .314 $\to$ .254; Phi: .275 $\to$ .251),
so the net effect of SECL is positive.
By comparison, supervised temperature scaling also reduces AUROC
(Llama: $-$.006; Gemma: $-$.008; Phi: $-$.017)
and Platt scaling degrades it further for Llama ($-$.029, below chance at .481),
demonstrating that AUROC loss is not unique to SECL but is an inherent
cost of recalibrating overconfident models.

\section{P(True) Normalization and the Generation--Discrimination Gap} \label{app:ptrue_norm}

A key premise of SECL is that LLMs possess a \emph{discriminative} 
confidence signal, P(True), that is better calibrated than their 
\emph{generative} verbalized confidence.
Table~\ref{tab:ptrue_raw_vs_norm} demonstrates this gap and the 
effect of normalization.

\paragraph{Raw P(True) vs.\ the verbalized baseline.}
For Llama, raw P(True) (the model's probability that its own answer 
is correct, without normalization) already improves calibration over 
the verbalized baseline (ECE: .161 vs.\ .170) and dramatically improves 
discrimination (AUROC: .673 vs.\ .510).
This confirms that the model "knows more than it says"; its 
internal judgment of correctness is more reliable than its 
expressed confidence.

\paragraph{Normalization with distractors.}
P(True) Norm applies a temperature-scaled softmax over the correct 
answer and $k{=}4$ distractors, converting a raw probability into a 
calibrated confidence score.
This yields a large further improvement (Llama ECE: .161 $\to$ .065 
at $\tau{=}0.7$), demonstrating that relative comparison against 
plausible alternatives sharpens the discriminative signal.

\begin{table}[htbp]
\centering
\small
\resizebox{\columnwidth}{!}{%
\begin{tabular}{lcccc}
\toprule
\textbf{Method}
  & \textbf{ECE}$\downarrow$ & \textbf{Brier}$\downarrow$
  & \textbf{AUROC}$\uparrow$ & \textbf{Acc} \\
\midrule
Verbalized        & .170 & .292 & .510 & .576 \\
Raw P(True)            & .161 & .248 & .673 & .576 \\
P(True) Norm $\tau{=}0.3$ & .118 & .238 & .685 & .576 \\
P(True) Norm $\tau{=}0.7$ & \textbf{.065} & \textbf{.223} & \textbf{.694} & .576 \\
P(True) Norm $\tau{=}1.5$ & .169 & .254 & .658 & .579 \\
\bottomrule
\end{tabular}}%
\caption{Effect of P(True) normalization (Llama~3.2-3B, $N{=}2{,}000$).
Raw P(True) improves over the verbalized baseline. Normalization with 
temperature $\tau$ yields further gains.
The optimal $\tau$ varies by model (see Table~\ref{tab:tau_sweep}).}
\label{tab:ptrue_raw_vs_norm}
\end{table}

\paragraph{Temperature selection.}
The normalization temperature $\tau$ controls sharpening ($\tau{<}1$) 
or smoothing ($\tau{>}1$) of the softmax over candidates.
Table~\ref{tab:tau_sweep} shows the P(True) Norm ECE at different 
$\tau$ values for each model.
Llama~3.2-3B benefits from sharpening ($\tau{=}0.7$), Gemma and Phi 
require smoothing ($\tau{=}1.5$), and Llama~3.1-8B requires stronger
smoothing ($\tau{=}3.0$), likely reflecting differences in 
the models' internal confidence distributions.
The selected $\tau$ is used in all subsequent experiments.

\begin{table}[htbp]
\centering
\small
\resizebox{\columnwidth}{!}{%
\begin{tabular}{lcccc}
\toprule
$\boldsymbol{\tau}$ & \textbf{Llama 3.2-3B} & \textbf{Gemma 2-2B} & \textbf{Phi 3.5-Mini} & \textbf{Llama 3.1-8B} \\
\midrule
0.3  & .118 & .255 & .172 & .270 \\
0.7  & \textbf{.065} & .215 & .168 & .234 \\
1.0  & .083 & .180 & .161 & .195 \\
1.5  & .169 & \textbf{.141} & \textbf{.154} & .153 \\
3.0  & .247 & .143 & .245 & \textbf{.120} \\
\bottomrule
\end{tabular}}%
\caption{P(True) Norm ECE as a function of normalization 
temperature $\tau$. Bold: selected value per model.}
\label{tab:tau_sweep}
\end{table}

\section{Hyperparameter Sensitivity} \label{app:hp_sensitivity}

Table~\ref{tab:hp_sensitivity} reports ECE and Adaptive ECE on 
Llama~3.2-3B when varying individual hyperparameters.
The $\alpha_{\text{step}}$ and $\delta$ ablations use the full pipeline 
\emph{without} entropy gating (all questions trained) to isolate the loss 
function's effect; the burst size $B$ ablation uses the entropy-gated pipeline.

\begin{table}[htbp]
\centering
\small
\resizebox{\columnwidth}{!}{%
\begin{tabular}{llcc}
\toprule
\textbf{Parameter} & \textbf{Value} & \textbf{ECE}$\downarrow$ & \textbf{AdaECE}$\downarrow$ \\
\midrule
\multirow{3}{*}{$\alpha_{\text{step}}$$^\dagger$}
  & 0.2 & .067 & .083 \\
  & 0.3 & .066 & .069 \\
  & \textbf{0.5} & \textbf{.052} & \textbf{.073} \\
\midrule
\multirow{2}{*}{$\delta$$^\dagger$}
  & \textbf{0.15} & \textbf{.052} & \textbf{.073} \\
  & 0.20 & .064 & .067 \\
\midrule
\multirow{1}{*}{Loss$^\dagger$}
  & Plain MSE (no directional) & .085 & .088 \\
\midrule
\multirow{2}{*}{Burst $B$$^\ddagger$}
  & 20 & .114 & -- \\
  & \textbf{50} & \textbf{.050} & \textbf{.060} \\
\midrule
\multirow{3}{*}{$\tau$ (per model)}
  & Llama: \textbf{0.7} & \textbf{.065} & \textbf{.064} \\
  & Gemma: \textbf{1.5} & .141 & .144 \\
  & Phi: \textbf{1.5} & .154 & .151 \\
\bottomrule
\end{tabular}}%
\caption{Hyperparameter sensitivity (Llama~3.2-3B).
Bold indicates the selected configuration.
$^\dagger$Without entropy gating.
$^\ddagger$With entropy gating.}
\label{tab:hp_sensitivity}
\end{table}

In additional non-default sweeps, we observed further gains for some models
(e.g., Llama in Table~\ref{tab:layer_ablation} and Phi under a tuned
$\lambda$/learning-rate pair, best observed ECE $0.097$). We keep the canonical
defaults in all headline comparisons to preserve a single cross-model protocol
and avoid per-model post-hoc tuning.

\newpage
\onecolumn
\section{Per-Domain Results} \label{app:per_domain}

\begin{table*}[ht]
\centering
\footnotesize
\begin{tabular}{lllccccc}
\toprule
\textbf{Model} & \textbf{Method} & \textbf{Domain} & \textbf{ECE}$\downarrow$ & \textbf{AdaECE}$\downarrow$ & \textbf{Brier}$\downarrow$ & \textbf{AUROC}$\uparrow$ & \textbf{Acc} \\
\midrule
\multirow{12}{*}{Llama 3.2-3B} & \multirow{4}{*}{Verbalized} 
  & GSM8K      & .218 & .215 & .294 & .565 & .472 \\
& & MMLU       & .106 & .109 & .253 & .601 & .562 \\
& & ARC        & \textbf{.095} & .112 & .207 & .568 & .726 \\
& & TruthfulQA & .372 & .371 & .414 & .301 & .544 \\
\cmidrule{2-8}
& \multirow{4}{*}{P(True) Norm} 
  & GSM8K      & .133 & .138 & .264 & .594 & .476 \\
& & MMLU       & .091 & .090 & .239 & .652 & .566 \\
& & ARC        & .117 & .132 & .210 & .624 & .728 \\
& & TruthfulQA & .089 & \textbf{.092} & .180 & .818 & .532 \\
\cmidrule{2-8}
& \multirow{4}{*}{SECL} 
  & GSM8K      & \textbf{.070} & \textbf{.079} & .249 & .573 & .488 \\
& & MMLU       & \textbf{.067} & \textbf{.060} & .249 & .549 & .558 \\
& & ARC        & .112 & \textbf{.106} & .207 & .616 & .714 \\
& & TruthfulQA & \textbf{.068} & .114 & .260 & .454 & .546 \\
\midrule
\multirow{12}{*}{Gemma 2-2B} & \multirow{4}{*}{Verbalized} 
  & GSM8K      & .549 & .549 & .446 & .658 & .186 \\
& & MMLU       & .194 & .178 & .271 & .608 & .576 \\
& & ARC        & \textbf{.082} & \textbf{.085} & .201 & .534 & .738 \\
& & TruthfulQA & .267 & .263 & .338 & .417 & .566 \\
\cmidrule{2-8}
& \multirow{4}{*}{P(True) Norm} 
  & GSM8K      & .395 & .394 & .326 & .685 & .186 \\
& & MMLU       & .163 & .175 & .263 & .616 & .576 \\
& & ARC        & .185 & .182 & .229 & .597 & .738 \\
& & TruthfulQA & \textbf{.104} & \textbf{.104} & .213 & .729 & .566 \\
\cmidrule{2-8}
& \multirow{4}{*}{SECL} 
  & GSM8K      & \textbf{.356} & \textbf{.356} & .271 & .609 & .180 \\
& & MMLU       & \textbf{.054} & \textbf{.064} & .238 & .594 & .578 \\
& & ARC        & .144 & .136 & .213 & .591 & .728 \\
& & TruthfulQA & .210 & .198 & .293 & .343 & .574 \\
\midrule
\multirow{12}{*}{Phi 3.5-Mini} & \multirow{4}{*}{Verbalized} 
  & GSM8K      & .290 & .290 & .304 & .563 & .650 \\
& & MMLU       & .264 & .257 & .286 & .602 & .648 \\
& & ARC        & \textbf{.109} & \textbf{.105} & .150 & .588 & .826 \\
& & TruthfulQA & .343 & .336 & .362 & .592 & .546 \\
\cmidrule{2-8}
& \multirow{4}{*}{P(True) Norm} 
  & GSM8K      & \textbf{.261} & \textbf{.257} & .303 & .559 & .650 \\
& & MMLU       & .171 & .158 & .243 & .648 & .648 \\
& & ARC        & .129 & .119 & .157 & .680 & .826 \\
& & TruthfulQA & \textbf{.146} & \textbf{.141} & .205 & .771 & .546 \\
\cmidrule{2-8}
& \multirow{4}{*}{SECL} 
  & GSM8K      & .283 & .280 & .297 & .575 & .658 \\
& & MMLU       & \textbf{.054} & \textbf{.058} & .223 & .604 & .644 \\
& & ARC        & .129 & .129 & .167 & .499 & .826 \\
& & TruthfulQA & .229 & .224 & .319 & .407 & .532 \\
\bottomrule
\end{tabular}
\caption{Per-domain results for Llama 3.2-3B, Gemma 2-2B, and Phi 3.5-Mini. Best ECE and AdaECE per domain are in \textbf{bold}.}
\label{tab:per_domain_combined}
\end{table*}

\clearpage

\twocolumn
\section{Ablation Studies} \label{app:ablations}

\paragraph{Weight accumulation.}
By default, SECL accumulates LoRA weights across the stream
(each question's adaptation builds on the previous state).
Table~\ref{tab:ablation_accum} shows that disabling accumulation
(resetting LoRA to the base model after each question)
is catastrophic: ECE degrades from 0.050 to 0.237, and AUROC
drops below chance (0.484), i.e., isolated per-question updates destroy the model's confidence
signal.

\begin{table}[htbp]
\centering
\small
\resizebox{\columnwidth}{!}{%
\begin{tabular}{lcccc}
\toprule
\textbf{Setting}
  & \textbf{ECE}$\downarrow$ & \textbf{Brier}$\downarrow$
  & \textbf{AUROC}$\uparrow$ & \textbf{Acc} \\
\midrule
SECL (accumulate)  & \textbf{.050} & \textbf{.241} & \textbf{.587} & .577 \\
Reset per question & .237 & .327 & .484 & .576 \\
\bottomrule
\end{tabular}}%
\caption{Weight accumulation ablation (Llama~3.2-3B).}
\label{tab:ablation_accum}
\end{table}

\paragraph{Entropy gating vs.\ always-train.}
Training on every question (no entropy gating) achieves comparable ECE
to SECL (0.047 vs.\ 0.050), but at 4$\times$ the computational cost
since all 2{,}000 questions receive TTT updates.
Entropy gating triggers calibration bursts when distribution shift is detected;
under our protocol this results in 25.6\% of questions receiving TTT updates,
achieving the same calibration with 75\% fewer gradient steps.

\begin{table}[htbp]
\centering
\small
\resizebox{\columnwidth}{!}{%
\begin{tabular}{lcccc}
\toprule
\textbf{Setting}
  & \textbf{ECE}$\downarrow$ & \textbf{Brier}$\downarrow$
  & \textbf{AUROC}$\uparrow$ & \textbf{Trained \%} \\
\midrule
Always-train   & \textbf{.047} & \textbf{.240} & \textbf{.593} & 100\% \\
Entropy-gated (SECL) & .050 & .241 & .587 & 25.6\% \\
\bottomrule
\end{tabular}}%
\caption{Entropy gating vs.\ always-train (Llama~3.2-3B).
Entropy gating matches always-train ECE at a fraction of the cost.}
\label{tab:ablation_phgate}
\end{table}

\paragraph{KL divergence regularization.}
We experiment with adding a KL divergence term
$\beta \cdot D_{\text{KL}}(p_{\text{base}} \| p_{\text{adapted}})$
to the training loss to preserve the base model's output distribution.
Table~\ref{tab:ablation_kl} shows that a small $\beta{=}0.01$ slightly
improves Llama ECE (0.044 vs.\ 0.050) but degrades Gemma and Phi,
while $\beta{=}0.1$ is too aggressive for all models.
We use $\beta{=}0$ (no KL term) in all reported experiments.

\begin{table}[htbp]
\centering
\small
\resizebox{\columnwidth}{!}{%
\begin{tabular}{llcccc}
\toprule
\textbf{Model} & $\boldsymbol{\beta}$
  & \textbf{ECE}$\downarrow$ & \textbf{Brier}$\downarrow$
  & \textbf{AUROC}$\uparrow$ & \textbf{Acc} \\
\midrule
\multirow{3}{*}{Llama 3.2-3B}
  & \textbf{0} (default) & .050 & \textbf{.241} & .587 & \textbf{.577} \\
  & 0.01     & \textbf{.044} & .242 & \textbf{.591} & .573 \\
  & 0.1      & .149 & .279 & .527 & .569 \\
\midrule
\multirow{3}{*}{Gemma 2-2B}
  & \textbf{0} (default) & \textbf{.056} & \textbf{.254} & .548 & .515 \\
  & 0.01     & .127 & .271 & .486 & .514 \\
  & 0.1      & .238 & .305 & \textbf{.551} & \textbf{.520} \\
\midrule
\multirow{3}{*}{Phi 3.5-Mini (3.8B)}
  & \textbf{0} (default) & \textbf{.110} & \textbf{.251} & .521 & .665 \\
  & 0.01     & .144 & .252 & .506 & .669 \\
  & 0.1      & .248 & .272 & \textbf{.598} & \textbf{.673} \\
\bottomrule
\end{tabular}}%
\caption{KL regularization ablation across models (SECL + entropy gating).}
\label{tab:ablation_kl}
\end{table}

\section{Ablation: LoRA Layer Position} \label{app:layer_ablation}

To justify placing LoRA adapters on late transformer layers,
we compare five layer configurations on Llama~3.2-3B while holding
all other hyperparameters fixed (top block of Table~\ref{tab:layer_ablation}).
Mid and late layers achieve the same ECE (0.039), while early layers
remain competitive (0.046).
Adapting all 28 layers yields the best AUROC and Brier
but the worst ECE (0.058), suggesting that adapting too many layers
introduces noise that hurts bin-level calibration despite improving
per-question discrimination.
The robustness across layer positions
indicates that SECL's calibration mechanism is not tightly coupled
to a specific layer group; we use late layers as the default following
evidence that calibration-relevant representations concentrate in
intermediate-to-late layers \citep{du_haloscope_2024}.

We further verify the effect of layer \emph{count} on Gemma~2-2B and
Phi~3.5-Mini (3.8B) (bottom block of Table~\ref{tab:layer_ablation}).
Halving the adapter count from 8 to 4 layers roughly doubles ECE on
both models (Gemma: 0.056$\to$0.116; Phi: 0.110$\to$0.222) with
negligible accuracy change, so 8 late layers provide
sufficient capacity for calibration across architectures.

\begin{table}[htbp]
\centering
\small
\resizebox{\columnwidth}{!}{%
\begin{tabular}{llccccc}
\toprule
\textbf{Model} & \textbf{Layers} & \textbf{Range}
  & \textbf{ECE}$\downarrow$ & \textbf{Brier}$\downarrow$
  & \textbf{AUROC}$\uparrow$ & \textbf{Acc} \\
\midrule
\multirow{5}{*}{Llama 3.2-3B}
  & Early     & 0--3   & .046 & .235 & .619 & .591 \\
  & Mid       & 12--15 & \textbf{.039} & .233 & .628 & .573 \\
  & Late (default) & 24--27 & \textbf{.039} & .241 & .592 & .570 \\
  & Last 8    & 20--27 & .044 & .240 & .609 & .558 \\
  & All       & 0--27  & .058 & \textbf{.232} & \textbf{.640} & \textbf{.588} \\
\midrule
\multirow{2}{*}{Gemma 2-2B}
  & Last 4    & 22--25 & .116 & .277 & .464 & .513 \\
  & Last 8 (default) & 18--25 & \textbf{.056} & \textbf{.254} & \textbf{.548} & \textbf{.515} \\
\midrule
\multirow{2}{*}{Phi 3.5-Mini (3.8B)}
  & Last 4    & 28--31 & .222 & .277 & \textbf{.524} & \textbf{.669} \\
  & Last 8 (default) & 24--31 & \textbf{.110} & \textbf{.251} & .521 & .665 \\
\bottomrule
\end{tabular}}%
\caption{LoRA layer ablation.
\textbf{Top:} layer position on Llama~3.2-3B (4 layers each except "All").
\textbf{Bottom:} layer count on Gemma~2-2B and Phi~3.5-Mini (3.8B).}
\label{tab:layer_ablation}
\end{table}

\section{Reliability Diagrams and Normalization Effect} \label{app:reliability}

Figure~\ref{fig:ptrue_norm} compares reliability diagrams for raw
$P(\text{True})$ and distractor-normalized
$\text{NormP}_{\text{True}}$ on Llama~3.2-3B.
Raw $P(\text{True})$ exhibits upward bias in mid-to-high confidence
bins, consistent with the suggestibility effect documented by
\citet{wang_calibrating_2025}: the model tends to affirm
plausible-sounding answers regardless of correctness.
Distractor normalization removes this bias by converting
absolute affirmation into relative preference among candidate
answers, reducing ECE from 0.161 to 0.065.

\begin{figure*}[htbp]
    \centering
    \includegraphics[width=0.85\textwidth]{figures/2_ptrue_vs_norm_llama.pdf}
    \caption{Raw $P(\text{True})$ (\textbf{left}) vs.\
    distractor-normalized $\text{NormP}_{\text{True}}$
    (\textbf{right}) on Llama~3.2-3B.
    Normalization suppresses suggestibility bias, producing a
    better-calibrated supervision signal
    (ECE: 0.161\,$\to$\,0.065).}
    \label{fig:ptrue_norm}
\end{figure*}

Figure~\ref{fig:ptrue_norm_all} extends this comparison to all three
models. Gemma~2-2B benefits the most (ECE: 0.369\,$\to$\,0.141),
while Phi~3.5-Mini (3.8B) shows a similarly large reduction
(ECE: 0.248\,$\to$\,0.154).

\begin{figure*}[htbp]
    \centering
    \includegraphics[width=0.85\textwidth]{figures/1_ptrue_vs_norm_all.pdf}
    \caption{Raw $P(\text{True})$ (\textbf{left}) vs.\ $\text{NormP}_{\text{True}}$ (\textbf{right}) for the three non-8B models. Distractor normalization consistently reduces calibration error across model families.}
    \label{fig:ptrue_norm_all}
\end{figure*}

Reliability diagrams for the verbalized baseline vs.\ SECL on all three
models are shown in Figure~\ref{fig:rel_all}.

\begin{figure*}[htbp]
    \centering
    \includegraphics[width=0.85\textwidth]{figures/4_secl_vs_verb_all.pdf}
    \caption{Reliability diagrams: Verbalized baseline (\textbf{left}) vs.\ SECL (\textbf{right}) for Llama~3.2-3B (ECE: 0.170\,$\to$\,0.050), Gemma~2-2B (ECE: 0.256\,$\to$\,0.056), and Phi~3.5-Mini (3.8B) (ECE: 0.251\,$\to$\,0.110). SECL consistently shifts predictions toward the diagonal, eliminating the systematic overconfidence visible in the verbalized baselines.}
    \label{fig:rel_all}
\end{figure*}

\section{Domain Order Sensitivity} \label{app:domain_order}

To assess sensitivity to the presentation order of domains, we run SECL 
with the reversed sequence 
(TruthfulQA $\to$ ARC $\to$ MMLU $\to$ GSM8K).
Table~\ref{tab:domain_order_full} compares overall metrics; 
Table~\ref{tab:domain_order_perdomain} gives per-domain detail.
SECL improves over the verbalized baseline under both orderings for all models,
though the original order yields lower ECE.
The gap is smallest for Llama~3.2-3B (0.050 vs.\ 0.068); Llama~3.1-8B shows
a larger gap (0.083 vs.\ 0.189), partly because the reversed order triggers
gating on 62.6\% of questions (vs.\ 12.6\% forward), amplifying cumulative
LoRA drift on the longer model.

\begin{table}[htbp]
\centering
\small
\resizebox{\columnwidth}{!}{%
\begin{tabular}{llccccc}
\toprule
\textbf{Model} & \textbf{Order}
  & \textbf{ECE}$\downarrow$ & \textbf{AdaECE}$\downarrow$ & \textbf{Brier}$\downarrow$
  & \textbf{AUROC}$\uparrow$ & \textbf{Acc} \\
\midrule
\multirow{3}{*}{Llama 3.2-3B}
  & Verbalized       & .170 & .167 & .292 & .510 & .576 \\
  & SECL $\mathcal{O}$ & \textbf{.050} & \textbf{.060} & .241 & .587 & .577 \\
  & SECL $\mathcal{R}$ & .068 & .082 & .248 & .578 & .575 \\
\midrule
\multirow{3}{*}{Gemma 2-2B}
  & Verbalized       & .256 & .252 & .314 & .558 & .516 \\
  & SECL $\mathcal{O}$ & \textbf{.056} & \textbf{.060} & .254 & .548 & .515 \\
  & SECL $\mathcal{R}$ & .149 & .133 & .258 & .625 & .520 \\
\midrule
\multirow{3}{*}{Phi 3.5-Mini (3.8B)}
  & Verbalized       & .251 & .240 & .275 & .600 & .667 \\
  & SECL $\mathcal{O}$ & \textbf{.110} & \textbf{.119} & .251 & .521 & .665 \\
  & SECL $\mathcal{R}$ & .173 & .170 & .250 & .575 & .674 \\
\midrule
\multirow{3}{*}{Llama 3.1-8B}
  & Verbalized       & .225 & .225 & .258 & .684 & .644 \\
  & SECL $\mathcal{O}$ & \textbf{.083} & \textbf{.080} & .222 & .643 & .646 \\
  & SECL $\mathcal{R}$ & .189 & .191 & .272 & .580 & .629 \\
\bottomrule
\end{tabular}}%
\caption{Domain order sensitivity (overall metrics).
$\mathcal{O}$: GSM8K$\to$MMLU$\to$ARC$\to$TruthfulQA.
$\mathcal{R}$: TruthfulQA$\to$ARC$\to$MMLU$\to$GSM8K.}
\label{tab:domain_order_full}
\end{table}

\begin{table}[htbp]
\centering
\small
\resizebox{\columnwidth}{!}{%
\begin{tabular}{llcccc}
\toprule
\textbf{Model} & \textbf{Order}
  & \textbf{GSM8K} & \textbf{MMLU}
  & \textbf{ARC} & \textbf{TruthfulQA} \\
\midrule
\multirow{2}{*}{Llama 3.2-3B}
  & $\mathcal{O}$ & .070 & .067 & .112 & .068 \\
  & $\mathcal{R}$ & .117 & .032 & .081 & .134 \\
\midrule
\multirow{2}{*}{Gemma 2-2B}
  & $\mathcal{O}$ & .356 & .054 & .144 & .210 \\
  & $\mathcal{R}$ & .396 & .096 & .046 & .189 \\
\midrule
\multirow{2}{*}{Phi 3.5-Mini (3.8B)}
  & $\mathcal{O}$ & .283 & .054 & .129 & .229 \\
  & $\mathcal{R}$ & .068 & .213 & .099 & .343 \\
\midrule
\multirow{2}{*}{Llama 3.1-8B}
  & $\mathcal{O}$ & .066 & .120 & .198 & .038 \\
  & $\mathcal{R}$ & .206 & .217 & .234 & .223 \\
\bottomrule
\end{tabular}}%
\caption{Per-domain ECE for original ($\mathcal{O}$) vs.\ reversed 
($\mathcal{R}$) domain order.}
\label{tab:domain_order_perdomain}
\end{table}

\paragraph{ARC domain ordering effect.}
ARC-Challenge is the only domain where SECL consistently increases ECE
relative to the verbalized baseline in the forward order
(Llama~3B: $.095 \to .112$; Gemma: $.082 \to .144$; Phi: $.109 \to .129$; Llama~8B: $.108 \to .198$).
In the reversed order, ARC ECE \emph{improves} for the three smaller models
(Llama~3B: $.095 \to .081$; Gemma: $.082 \to .046$; Phi: $.109 \to .099$),
though Llama~8B degrades further ($.108 \to .234$).
For the smaller models, the ARC degradation is therefore an artifact of domain sequencing, specifically,
the LoRA weights adapted on preceding domains transfer poorly to ARC's
scientific reasoning format, not a limitation of SECL itself.
Overall ECE remains improved under both orderings for all models,
so the per-domain redistribution does not undermine the method's aggregate benefit.

\section{Computational Cost Analysis} \label{app:cost}

Table~\ref{tab:cost} compares the computational cost of SECL against
running P(True) Norm on every question.
We report costs in forward-pass equivalents (FWD-eq), where one 
backward pass $\approx$ 2 FWD.
For questions where entropy gating skips TTT, SECL requires only 1~FWD 
(the generation pass); the P(True) Norm baseline requires $1 + (k{+}1) = 6$~FWD 
per question ($k{=}4$ distractors plus the correct answer).

\begin{table}[htbp]
\centering
\small
\resizebox{\columnwidth}{!}{%
\begin{tabular}{lcccc}
\toprule
\textbf{Model} & \textbf{Trained} & \textbf{Skipped}
  & \textbf{SECL} & \textbf{P(True)} \\
\midrule
Llama 3.2-3B & 512 (25.6\%) & 1{,}488 & 9{,}168 & 12{,}000 \\
Llama 3.1-8B & 251 (12.6\%) & 1{,}749 & 5{,}514 & 12{,}000 \\
Gemma 2-2B   & 119 ~~(5.9\%) & 1{,}881 & 3{,}666 & 12{,}000 \\
Phi 3.5-Mini (3.8B) & 160 ~~(8.0\%) & 1{,}840 & 4{,}240 & 12{,}000 \\
\bottomrule
\end{tabular}}%
\caption{Computational cost comparison (in FWD-equivalents).
\emph{Trained} = questions receiving TTT;
\emph{Skipped} = entropy-gated questions (generation only).
Per trained question: ${\approx}15$ FWD-eq 
(1~gen $+$ 5~P(True) $+$ 3~epochs $\times$ 3~FWD-eq/step).
Per skipped question: $1$ FWD-eq.
P(True) Norm baseline: 6 FWD-eq $\times$ 2{,}000 $= 12{,}000$.}
\label{tab:cost}
\vspace{-0.5em}
\end{table}

SECL is cheaper than P(True) Norm baseline for all four models.
Gemma and Phi trigger infrequently (6--8\%), reducing cost to roughly
a third of the baseline while achieving substantially better calibration.
Even Llama, which triggers most often (25.6\%), remains below the
baseline cost while improving ECE from .065 to .050.
The temperature scaling and Platt scaling add negligible compute
(a single parameter fit) but require ground-truth labels, which are 
unavailable in the test-time setting.

\section{Negative Control: Qwen 2.5-3B} \label{app:qwen}

We evaluate Qwen~2.5-3B-Instruct \citep{qwen_qwen25_2025} as a negative control.
Unlike the other four models, Qwen's P(True) Norm baseline 
(\emph{best} $\tau{=}1.0$, ECE$=$0.257) is \emph{worse} than its
verbalized baseline (ECE$=$0.247), indicating that the 
discriminative signal provides no calibration advantage over generation.
Since SECL distills the discriminative signal into generative confidence,
the method cannot improve calibration when the discriminative signal 
itself is uninformative.
This confirms a key prerequisite of our approach: SECL succeeds 
precisely when a generation--discrimination gap exists.

\begin{table}[htbp]
\centering
\small
\begin{tabular}{lc}
\toprule
\textbf{Method} & \textbf{ECE}$\downarrow$ \\
\midrule
Verbalized & \textbf{.247} \\
\addlinespace
P(True) Norm $\tau{=}0.3$ & .290 \\
P(True) Norm $\tau{=}0.7$ & .263 \\
P(True) Norm $\tau{=}1.0$ & .257 \\
P(True) Norm $\tau{=}1.5$ & .265 \\
P(True) Norm $\tau{=}2.0$ & .267 \\
P(True) Norm $\tau{=}3.0$ & .291 \\
\bottomrule
\end{tabular}
\caption{Qwen~2.5-3B: no generation--discrimination gap.
All P(True) Norm temperatures yield worse ECE than the verbalized baseline.}
\label{tab:qwen}
\end{table}

\clearpage

\section{Entropy Gating Trigger Statistics} \label{app:ph_stats}

\begin{table}[htbp]
\centering
\small
\resizebox{\columnwidth}{!}{%
\begin{tabular}{lccc}
\toprule
\textbf{Model} & \textbf{Total} & \textbf{Trained (\%)} & \textbf{Skipped (\%)} \\
\midrule
Llama 3.2-3B & 2{,}000 & 512 (25.6\%) & 1{,}488 (74.4\%) \\
Gemma 2-2B   & 2{,}000 & 119 ~~(5.9\%) & 1{,}881 (94.1\%) \\
Phi 3.5-Mini (3.8B) & 2{,}000 & 160 ~~(8.0\%) & 1{,}840 (92.0\%) \\
Llama 3.1-8B & 2{,}000 & 251 (12.6\%) & 1{,}749 (87.4\%) \\
\bottomrule
\end{tabular}}%
\caption{Entropy gating trigger statistics per model.
\emph{Trained\%} is the fraction of questions receiving TTT adaptation.}
\label{tab:ph_triggers}
\end{table}

\section{Scaling to Llama 3.1-8B} \label{app:scaling_8b}

To test whether SECL scales to larger models, we evaluate
Llama~3.1-8B-Instruct (8.03B parameters) on the same 2{,}000-question stream.
We use LoRA on the last 8 layers (rank~8), $\tau{=}3.0$,
and the same entropy gating parameters as the 3B models.

\begin{table}[htbp]
\centering
\small
\resizebox{\columnwidth}{!}{%
\begin{tabular}{lccccc}
\toprule
\textbf{Method}
  & \textbf{ECE}$\downarrow$ & \textbf{AdaECE}$\downarrow$ & \textbf{Brier}$\downarrow$
  & \textbf{AUROC}$\uparrow$ & \textbf{Acc} \\
\midrule
Verbalized      & .225 & .225 & .258 & .684 & .644 \\
~~+ Temp Scaling$^*$ & \textbf{.080} & .089 & .213 & .681 & .644 \\
P(True) Norm         & .120 & .117 & .211 & \textbf{.718} & .646 \\
~~+ Temp Scaling$^*$ & .108 & .103 & \textbf{.208} & .716 & .646 \\
SECL (Ours)          & \underline{.083} & \textbf{.080} & .222 & .643 & .646 \\
\bottomrule
\end{tabular}}%
\caption{Llama~3.1-8B results. SECL reduces ECE by 63\% relative
to the verbalized baseline while adapting only 12.6\% of questions.
$^*$Requires ground-truth labels (5-fold CV).}
\label{tab:8b_results}
\end{table}

\begin{table}[htbp]
\centering
\small
\resizebox{\columnwidth}{!}{%
\begin{tabular}{llcccc}
\toprule
\textbf{Method} & \textbf{Domain}
  & \textbf{ECE}$\downarrow$ & \textbf{Brier}$\downarrow$
  & \textbf{AUROC}$\uparrow$ & \textbf{Acc} \\
\midrule
\multirow{4}{*}{Verbalized}
  & GSM8K      & .298 & .305 & .714 & .602 \\
  & MMLU       & .208 & .251 & .644 & .670 \\
  & ARC        & .108 & .164 & .633 & .798 \\
  & TruthfulQA & .287 & .310 & .669 & .506 \\
\midrule
\multirow{4}{*}{SECL}
  & GSM8K      & \textbf{.066} & .238 & .602 & .572 \\
  & MMLU       & \textbf{.120} & .223 & .591 & .686 \\
  & ARC        & .198 & .178 & .632 & .824 \\
  & TruthfulQA & \textbf{.038} & .249 & .552 & .504 \\
\bottomrule
\end{tabular}}%
\caption{Per-domain results for Llama~3.1-8B.}
\label{tab:8b_per_domain}
\end{table}

SECL achieves ECE$=$0.083 on the 8B model, a 63\% reduction from the
verbalized baseline (0.225), confirming that the method scales
effectively. The pattern mirrors the 3B results: large improvements
on GSM8K, MMLU, and TruthfulQA, with ARC remaining resistant to
improvement (see domain ordering analysis in Section~\ref{app:domain_order}).
The entropy gate triggers on 12.6\% of questions, intermediate between
the 3B models' 5.9--25.6\% range.

\section{Ablation: Training Target Signal}
\label{app:target_ablation}

A natural question is whether P(True) Norm is the right pseudo-label
for SECL training.
Self-Consistency \citep[SC;][]{wang_selfconsistency_2023} is an alternative unsupervised signal that
measures agreement among $N{=}10$ temperature-sampled answers.
We run the full SECL pipeline (LoRA adaptation, entropy gating,
burst training, MSE loss) with SC as the training target, keeping
all other hyperparameters identical to the canonical configuration.

\begin{table}[htbp]
\centering
\small
\resizebox{\columnwidth}{!}{%
\begin{tabular}{lcccc}
\toprule
\textbf{Training Target}
  & \textbf{ECE}$\downarrow$ & \textbf{Brier}$\downarrow$
  & \textbf{AUROC}$\uparrow$ & \textbf{Acc} \\
\midrule
None (Verbalized baseline)     & .170 & .292 & .510 & .576 \\
Self-Consistency ($N{=}10$) & .432 & .470 & .443 & .564 \\
P(True) Norm (ours)        & \textbf{.050} & \textbf{.241} & \textbf{.587} & \textbf{.577} \\
\bottomrule
\end{tabular}}%
\caption{Effect of training target on Llama~3.2-3B.
All rows use the same SECL pipeline; only the pseudo-label
differs. Self-Consistency as a target degrades calibration below
the untrained baseline.}
\label{tab:target_ablation}
\end{table}

\paragraph{Why the target signal determines calibration quality.}
SECL minimizes $\mathcal{L} = \mathbb{E}[(c(q) - t(q))^2]$, where
$c(q)$ is the model's verbalized confidence and $t(q)$ is the
training target.
With sufficient model capacity and training convergence,
$c^*(q) \to t(q)$, so the trained model \emph{inherits the
calibration of its target}:
$\mathrm{ECE}(c^*) \to \mathrm{ECE}(t)$.
The choice of target therefore bounds the achievable calibration.

\paragraph{SC is a biased proxy for correctness.}
As $N \to \infty$, SC$(q) \to p_{\mathrm{gen}}(\text{mode} \mid q)$,
the probability mass the generation distribution places on its modal
answer.
For modern LLMs, the generation distribution is low-entropy:
$p_{\mathrm{gen}}(\text{mode} \mid q)$ is typically high
($>$0.7) regardless of whether the answer is correct.
Among questions where SC~$\approx 1$, the actual fraction correct
can be significantly below~1, violating the calibration condition
$P(Y{=}1 \mid c{=}p) = p$.
Training on SC therefore teaches the model to report generation
concentration as confidence, producing systematic overconfidence.

\begin{figure*}[ht]
\centering
\includegraphics[width=0.85\textwidth]{figures/6_sc_vs_ptrue_vs_norm_llama.pdf}
\caption{Reliability diagrams comparing candidate training targets on Llama~3.2-3B. Self-Consistency (\textbf{left}, ECE$=$0.093) produces a reasonable calibration curve but with systematic bias. Raw P(True) (\textbf{center}, ECE$=$0.161) exhibits upward bias from suggestibility. P(True) Norm (\textbf{right}, ECE$=$0.065) tracks the diagonal most closely, confirming it as the best-calibrated supervision signal for SECL training.}
\label{fig:sc_ptrue_norm_llama}
\end{figure*}

\begin{figure*}[ht]
\centering
\includegraphics[width=0.85\textwidth]{figures/target_distributions.pdf}
\caption{Confidence score distributions for correct (green) and incorrect (red) predictions on Llama~3.2-3B. Self-Consistency and Raw P(True) both cluster near 1.0 with poor class separation; P(True) Norm spreads scores across $[0,1]$ with better separation between correct and incorrect predictions.}
\label{fig:target_distributions}
\end{figure*}

\paragraph{P(True) Norm is a less biased proxy.}
P(True) Norm computes $\mathrm{softmax}(\log P(\text{True} \mid
q, a_i) / \tau)_0$ over the model's answer and $k$ distractors.
When the model cannot distinguish its answer from distractors
(equal $P(\text{True})$ scores), P(True) Norm~$\to 1/(k{+}1)$,
correctly indicating low confidence.
When the model strongly prefers its answer,
P(True) Norm~$\to 1$.
This normalization maps the discriminative signal to $[0,1]$, aligning it more closely with the correctness probability and reducing both raw and trained ECE.

The raw P(True) Norm signal (ECE$=$0.065) is better calibrated than the
raw Self-Consistency signal (ECE$=$0.093) \emph{before any training}.
After SECL training with the P(True) Norm target, the model's verbalized
confidence inherits this calibration (ECE$=$0.050), while training
With the SC target, SC's biases are amplified (ECE$=$0.432).
Figure~\ref{fig:sc_ptrue_norm_llama} provides direct visual
evidence via reliability diagrams: SC produces a monotonically
increasing but biased calibration curve, Raw P(True) suffers from
suggestibility bias in mid-to-high bins, while P(True) Norm
tracks the diagonal most closely.
Figure~\ref{fig:target_distributions} complements this with
score distributions, showing that P(True) Norm achieves the best
separation between correct and incorrect predictions.

Calibration quality is thus determined by the training target.
P(True) Norm is a better proxy for $P(\text{correct})$ than Self-Consistency; the adaptation mechanism distills whatever signal it receives, so signal quality determines the ceiling of SECL's effectiveness.

\section{Extended Stream Length (4{,}000 Questions)}
\label{app:extended_stream}

To test whether SECL's calibration improvements are stable over longer streams, we double the stream length from 2{,}000 to 4{,}000 questions (1{,}000 per domain) on Llama~3.2-3B, Gemma~2-2B, and Phi~3.5-Mini (3.8B). Table~\ref{tab:extended_stream} reports the results.

\begin{table}[htbp]
\centering
\small
\resizebox{\columnwidth}{!}{%
\begin{tabular}{llcccc}
\toprule
\textbf{Model} & \textbf{Method} & \textbf{ECE}$\downarrow$ & \textbf{Brier}$\downarrow$ & \textbf{AUROC}$\uparrow$ & \textbf{Acc} \\
\midrule
\multirow{2}{*}{Llama 3.2-3B}
  & Verbalized (4{,}000q) & .190 & .313 & .508 & .577 \\
  & SECL (4{,}000q) & \textbf{.051} & \textbf{.239} & \textbf{.607} & \textbf{.581} \\
\midrule
\multirow{2}{*}{Gemma 2-2B}
  & Verbalized (4{,}000q) & .276 & .328 & \textbf{.548} & .508 \\
  & SECL (4{,}000q) & \textbf{.086} & \textbf{.263} & .500 & \textbf{.512} \\
\midrule
\multirow{2}{*}{Phi 3.5-Mini (3.8B)}
  & Verbalized (4{,}000q) & \textbf{.192} & \textbf{.245} & .593 & .687 \\
  & SECL (4{,}000q) & .259 & .260 & \textbf{.712} & \textbf{.696} \\
\bottomrule
\end{tabular}}%
\caption{4{,}000-question comparison against same-stream verbalized baselines. SECL substantially improves calibration on Llama and Gemma. For Phi, stronger long-stream calibration requires sufficient distillation; with the 4{,}000q run, calibration remains worse than verbalized while discrimination and accuracy improve.}
\label{tab:extended_stream}
\end{table}

\begin{table}[t]
\centering
\small
\resizebox{\columnwidth}{!}{%
\begin{tabular}{llccccc}
\toprule
\textbf{Model} & \textbf{Method} & \textbf{FWD-eq}
  & \textbf{ECE}$\downarrow$ & \textbf{Brier}$\downarrow$
  & \textbf{AUROC}$\uparrow$ & \textbf{Acc} \\
\midrule
\multirow{3}{*}{Llama 3.2-3B}
  & Verbalized & 1          & .170 & .292 & .510 & .576 \\
  & DINCO           & $\sim$10   & .101 & \textbf{.207} & \textbf{.762} & .466 \\
  & SECL (ours)     & 4.6        & \textbf{.050} & .241 & .587 & \textbf{.577} \\
\midrule
\multirow{3}{*}{Phi 3.5-Mini (3.8B)}
  & Verbalized & 1          & .251 & .275 & .600 & \textbf{.667} \\
  & DINCO           & $\sim$10   & \textbf{.110} & \textbf{.212} & \textbf{.749} & .560 \\
  & SECL (ours)     & 2.1        & \textbf{.110} & .251 & .521 & .665 \\
\midrule
\multirow{3}{*}{Gemma 2-2B}
  & Verbalized & 1          & .256 & .314 & .558 & \textbf{.516} \\
  & DINCO           & $\sim$10   & .408 & .410 & .566 & .327 \\
  & SECL (ours)     & 1.8        & \textbf{.056} & \textbf{.254} & .548 & .515 \\
\midrule
\multirow{3}{*}{Llama 3.1-8B}
  & Verbalized & 1          & .225 & .258 & .684 & .644 \\
  & DINCO           & $\sim$10   & .117 & \textbf{.210} & \textbf{.756} & .522 \\
  & SECL (ours)     & 2.8        & \textbf{.083} & .222 & .643 & \textbf{.646} \\
\bottomrule
\end{tabular}}%
\vspace{0.3em}
\par\noindent{\small \emph{FWD-eq}: amortized forward-pass equivalents per question over the full stream (see Section~\ref{app:cost}).}
\caption{SECL vs.\ DINCO across all models.
SECL achieves better calibration (ECE) at lower amortized cost,
while DINCO provides superior discrimination (AUROC) at higher cost.}
\label{tab:dinco}
\end{table}

At 4{,}000 questions, SECL improves over verbalized confidence on Llama (ECE: 0.190$\to$0.051; Brier: 0.313$\to$0.239; AUROC: 0.508$\to$0.607) and Gemma (ECE: 0.276$\to$0.086; Brier: 0.328$\to$0.263), though Gemma shows a modest AUROC trade-off (0.548$\to$0.500). For Phi, SECL improves AUROC/accuracy (0.593$\to$0.712; 0.687$\to$0.696) but remains worse on calibration (ECE/Brier: 0.259/0.260 vs.\ 0.192/0.245), indicating under-distillation.

\section{Additional Open-Ended Evaluation (TruthfulQA-Gen)} \label{app:tqgen}

To reduce the gap between multiple-choice evaluation and open-ended generation, we replace the final domain with the TruthfulQA generation split and keep the rest of the stream unchanged. We report the forward and reversed streams (500 questions per domain, 2{,}000 total) and an extended 4{,}000-question forward stream (1{,}000 per domain) to test long-horizon stability on open-ended answers.

\begin{table}[htbp]
\centering
\small
\resizebox{\columnwidth}{!}{%
\begin{tabular}{llccccc}
\toprule
\textbf{Method} & \textbf{Order} & \textbf{Stream} & \textbf{ECE}$\downarrow$ & \textbf{Brier}$\downarrow$ & \textbf{AUROC}$\uparrow$ & \textbf{Acc} \\
\midrule
Verbalized & Forward & 2k & .195 & .287 & .586 & \textbf{.506} \\
P(True) Norm & Forward & 2k & .173 & .273 & \textbf{.625} & \textbf{.506} \\
SECL (Ours) & Forward & 2k & \textbf{.047} & \textbf{.239} & .618 & .500 \\
SECL (Ours) & Reversed & 2k & \textbf{.047} & .241 & \textbf{.625} & .497 \\
\midrule
Verbalized & Forward & 4k & .179 & .278 & .598 & \textbf{.525} \\
SECL (Ours) & Forward & 4k & \textbf{.061} & \textbf{.237} & \textbf{.625} & .522 \\
\bottomrule
\end{tabular}}%
\caption{TruthfulQA generation-domain evaluation (Llama~3.2-3B).
SECL reduces ECE by 76\% on the 2k forward stream (0.195\,$\to$\,0.047), remains stable under reversed ordering, and at 4k still improves over the same-stream verbalized baseline by 66\% (0.179\,$\to$\,0.061).}
\label{tab:tqgen}
\end{table}

\section{Comparison with DINCO}
\label{app:dinco}

DINCO \citep{wang_calibrating_2025} is a recent inference-time calibration
method that normalizes verbalized confidence against self-generated
distractors and combines it with Self-Consistency.
Unlike SECL, DINCO does not adapt the model; it computes a calibrated
confidence score at each inference step using $\sim$10 forward passes
(beam search + P(True) per candidate + NLI + SC sampling).

SECL and DINCO address different trade-offs.
Across all four models, SECL achieves lower or equal ECE
(Llama~3B: 0.050 vs.\ 0.101; Llama~8B: 0.083 vs.\ 0.117;
Phi: 0.110 vs.\ 0.110; Gemma: 0.056 vs.\ 0.408)
with lower total computational cost and no external models.
DINCO achieves stronger AUROC on all four models
using multiple inference passes and an NLI model,
at the cost of $\sim$10$\times$ inference compute.
The calibration--cost trade-off is visualized in Figure~\ref{fig:pareto_main} (main paper).

Under DINCO's default configuration, accuracy is consistently lower than SECL's
(Llama~3B: 46.6\% vs.\ 57.7\%; Llama~8B: 52.2\% vs.\ 64.6\%;
Phi: 56.0\% vs.\ 66.5\%; Gemma: 32.7\% vs.\ 51.5\%),
as beam-search answer selection yields lower task
accuracy than greedy decoding; using greedy answers with beam-search distractors could mitigate this.
On Gemma, DINCO essentially fails: ECE of 0.408 is worse than
the verbalized baseline (0.256), and accuracy drops to 32.7\%.
Gemma's internal representations appear ill-suited
to DINCO's NLI-based reweighting, whereas SECL's distillation
of the discriminative signal still works (ECE 0.056).

